\chapter{Флуктуационный детерминант моста и отрицательная мода чистоты}

Классическое матричное действие моста стабилизирует изотропию. Однако
проект уже допускает другую, ранее построенную связь состояния с
оператором: семейное состояние может входить как нормированный вес следа.
Настоящая проверка выясняет, способен ли однопетлевой детерминант физических
флуктуаций моста изменить знак ориентационного взаимодействия.

\section{Состояние как вес одного родительского следа}

На пятимерной кратности и семейном триплете состояние
\begin{equation}
 \tau_5\otimes R,qquad R=R^T>0,\qquad\Tr R=1,
 \label{eq:v6-state-weighted-bridge-trace}
\end{equation}
является нормированным и при $R=I_3/3$ возвращает прежний равномерный
след. Без нового коэффициента матричное действие принимает вид
\begin{equation}
 \mathcal S_R(B)=\frac17\left[
 \Tr R(I-B^TB)^2+\Tr R(I-BB^T)^2\right].
 \label{eq:v6-state-weighted-bridge-action}
\end{equation}
При $B=\lambda I_3$ оно точно равно
$2(1-\lambda^2)^2/7$ для любого $R$.

Это отличается от условного композитного словаря $R=B^TB/3$ предыдущего
проверки. Здесь $R$ является независимым состоянием, оценивающим кривизну
моста, как в уже построенной вариации семейного состояния.

\section{Физический гессиан замкнутого моста}

Положим $B=I_3+H$. До второго порядка
\begin{equation}
 \mathcal S_R(I_3+H)
 =\frac27\Tr R(H+H^T)^2+O(H^3).
\end{equation}
Три антисимметричных направления касательны к минимумам $O(3)$ и являются
нулевыми модами. На шести симметричных физических модах, в базисе, где
$R=\operatorname{diag}(r_1,r_2,r_3)$, ненулевые собственные значения
гессиана равны
\begin{align}
 m_{ii}^2&=\frac{16}{7}r_i,
 &i&=1,2,3,\label{eq:v6-bridge-diagonal-masses}\\
 m_{ij}^2&=\frac87(r_i+r_j),
 &i&<j.\label{eq:v6-bridge-offdiagonal-masses}
\end{align}
При $R=I_3/3$ все шесть равны $16/21$, что совпадает с прямым гессианом
канонического моста после удаления трёх орбитальных нулей.

Однопетлевая поправка в плоской мере этих физических мод имеет вид
\begin{equation}
 \Gamma_1(R)=\frac12\log\det{}'\Hess_B\mathcal S_R
 =\frac12\left[
 \sum_i\log r_i+\sum_{i<j}\log(r_i+r_j)
 \right]+\text{const}.
 \label{eq:v6-bridge-one-loop-determinant}
\end{equation}
Это стандартная background-field структура
$\frac12\Tr\log\Hess$; см. \path{arXiv:1505.03119} и литературу по
интегрированию тяжёлых полей методом фонового поля.

\section{Точный знак около изотропного состояния}

Пусть
\begin{equation}
 R=\frac13I_3+\delta R,qquad\Tr\delta R=0.
\end{equation}
Разложение \eqref{eq:v6-bridge-one-loop-determinant} даёт
\begin{equation}
 \Gamma_1(R)-\Gamma_1(I_3/3)
 =-\frac{45}{16}\Tr(\delta R)^2+O(\delta R^3).
 \label{eq:v6-induced-negative-purity}
\end{equation}
Следовательно, индуцированный ориентационный коэффициент равен
\begin{equation}
 \boxed{\kappa_{\rm fluc}=\frac{45}{16}=2.8125}.
\end{equation}
Он превышает оба ранее полученных порога:
\begin{align}
 \frac{45}{16}&>\log4,\\
 \frac{45}{16}&>\log4+\frac67
 =2.243437218\ldots.
\end{align}

Энтропия даёт положительный квадратичный коэффициент $3/2$. Даже если
сохранить положительную композитную цену $6/7$, суммарный коэффициент
равен
\begin{equation}
 \frac32+\frac67-\frac{45}{16}
 =-\frac{51}{112}<0.
 \label{eq:v6-total-negative-purity-hessian}
\end{equation}

\begin{theorem}[Условная однопетлевая неустойчивость изотропного состояния]
При состоянии-взвешенном следе \eqref{eq:v6-state-weighted-bridge-trace},
плоской мере физических симметричных флуктуаций и стандартном коэффициенте
$\frac12\log\det{}'$, замкнутый мост индуцирует отрицательную моду чистоты,
достаточную для потери устойчивости $R=I_3/3$.
\end{theorem}

Это первый найденный внутри общей архитектуры механизм, который способен
запустить проекторное упорядочение без внешнего толчка и без явного выбора
оси.

\section{Почему рождение материи ещё не доказано}

Поправка \eqref{eq:v6-bridge-one-loop-determinant} стремится к
$-\infty$, когда одно из $r_i$ стремится к нулю. Одновременно
соответствующая масса \eqref{eq:v6-bridge-diagonal-masses} обращается в
нуль. Поэтому однопетлевое приближение само ломается именно там, куда
направляет состояние. Оно доказывает локальную неустойчивость изотропной
точки, но не конечный устойчивый одноосный минимум.

Кроме того, Том V не вывел полную полевую/BV-меру. Плоская мера на шести
симметричных модах является минимальным контрольным выбором, но изменение
меры или якобиан полярного разложения может добавить зависящие от $R$
члены. Абсолютный логарифм гессиана также требует общей размерной
нормировки, хотя коэффициент относительного разложения около $I_3/3$ от
постоянного масштаба не зависит.

\begin{nogocriterion}[Граница однопетлевого запуска]
Нельзя объявлять материю рождённой только по отрицательному гессиану
\eqref{eq:v6-total-negative-purity-hessian}. Требуется получить конечный
нелинейно насыщенный минимум полного матричного интеграла, вывести меру и
проверить, что его орбита равна $\mathbb{RP}^2$, после чего быстрое переключение оставляет
локальную пару $+15/-15$.
\end{nogocriterion}

\begin{protocol}
\begin{enumerate}
 \item состояние-взвешенный нормированный след:
 \textbf{совместим с проектной архитектурой};
 \item новый свободный классический коэффициент:
 \textbf{не введён};
 \item физические мостовые моды после $O(3)$-фактора:
 \textbf{шесть};
 \item их гессиан:
 \textbf{вычислен точно};
 \item однопетлевой коэффициент чистоты:
 \textbf{$-45/16$};
 \item превышение порога $\log4+6/7$:
 \textbf{да};
 \item локальная устойчивость $I_3/3$:
 \textbf{теряется условно};
 \item внешний временной толчок:
 \textbf{для локального запуска не требуется};
 \item конечный одноосный минимум:
 \textbf{не получен};
 \item полная мера и BV-фактор:
 \textbf{не выведены};
 \item следующая проверка:
 \textbf{непертурбативное насыщение состояния-взвешенного матричного моста}.
\end{enumerate}
\end{protocol}

Машинный сертификат сохранён в
\path{s2t_v6_bridge_fluctuation_determinant_purity_gate_results.json}.